\begingroup
\renewcommand{\arraystretch}{1.45}
\clearpage\fancyhead[L]{\small DETERMINATION OF ATTRACTION}\thispagestyle{empty}\pdfbookmark[0]{Determination of Attraction}{attraction-title}\SourceLeaf{340}
\vspace*{1.5em}
\begin{center}
\textbf{DETERMINATION OF THE ATTRACTION}

ON ANY POINT OF GIVEN POSITION

\textbf{THAT A PLANET WOULD EXERT IF ITS MASS}

AROUND ITS ENTIRE ORBIT

\textbf{IN PROPORTION TO THE TIME IN WHICH ITS INDIVIDUAL PARTS ARE TRAVERSED}

WERE UNIFORMLY DISTRIBUTED

\bigskip
BY

\bigskip
CARL FRIEDRICH GAUSS

\bigskip
PRESENTED TO THE ROYAL SOCIETY OF SCIENCES IN 1818, ON JANUARY 17.

\bigskip
\rule{0.7\linewidth}{0.4pt}

Commentationes societatis regiae scientiarum Gottingensis recentiores. Vol.~\textsc{iv}.

Göttingen \textsc{mdcccxviii}.

\rule{0.7\linewidth}{0.4pt}
\end{center}
\clearpage\fancyhead[L]{\small DETERMINATION OF ATTRACTION}\SourceLeaf{342}
{\small\begin{center}
DETERMINATION OF THE ATTRACTION\\
ON ANY POINT OF GIVEN POSITION\\
THAT A PLANET WOULD EXERT IF ITS MASS\\
AROUND ITS ENTIRE ORBIT\\
IN PROPORTION TO THE TIME IN WHICH ITS INDIVIDUAL PARTS ARE TRAVERSED\\
WERE UNIFORMLY DISTRIBUTED.
\end{center}}

\begin{center}\rule{2cm}{0.3pt}\end{center}
\Needspace{23\baselineskip}
\AttractionArticle{1}
The secular variations that the elements of a planetary orbit undergo through perturbation by another planet are independent of the latter’s position in its orbit, and would be the same whether the perturbing planet moved in an elliptical orbit according to \textsc{Kepler’s} laws or its mass were conceived as distributed uniformly around the orbit in such a way that equal parts of the mass were now assigned to parts of the orbit otherwise traversed in equal intervals of time, provided that the periods of revolution of the perturbed and perturbing planets are not commensurable. This elegant theorem, if no one has so far stated it in express words, can at least be demonstrated very easily from the principles of physical astronomy. A problem therefore presents itself that is well worth attention both in itself and because of the several devices its solution requires: to determine exactly the attraction of a planetary orbit, or, if you prefer. of an elliptical ring whose thickness is infinitely small and variable according to the law just explained, on any point of given position.

\Needspace{40\baselineskip}
\AttractionArticle{2}
With the eccentricity of the orbit denoted by $e$ and the eccentric anomaly of any point in it by $E$, the element $\mathrm{d}E$ of the latter will correspond to the element of mean anomaly $(1-e\cos E)\mathrm{d}E$; consequently, the element of mass to be assigned to the small portion of the orbit to which these elements correspond will be to the entire mass, which we shall take as the unit, as $(1-e\cos E)\mathrm{d}E$ is to $2\pi$, where $\pi$ expresses the semicircumfe\SourceLeaf{343}%
rence of a circle of radius 1. Thus, setting the distance of the attracted point from the point of the orbit $=\rho$, the attraction produced by the orbital element will be
\[
\begin{aligned}
&=\frac{(1-e\cos E)\mathrm{d}E}{2\pi\rho\rho}
\end{aligned}
\]
We shall denote the major semiaxis by $a$ and the minor semiaxis by $b$, and adopt the former as the line of abscissas and the centre of the ellipse as the origin. Hence $aa-bb=aaee$ will hold, with the abscissa of the point of the orbit $=a\cos E$ and its ordinate $=b\sin E$. Finally, we shall denote the distance of the attracted point from the plane of the orbit by $C$, and its remaining coordinates parallel to the major and minor axes by $A$ and $B$. With these preparations, the attraction of the orbital element will be resolved into two attractions parallel to the major and minor axes and a third normal to the plane of the orbit, namely
\[
\begin{aligned}
&\frac{(A-a\cos E)(1-e\cos E)\mathrm{d}E}{2\pi\rho^3}=\mathrm{d}\xi\\[6pt]
&\frac{(B-b\sin E)(1-e\cos E)\mathrm{d}E}{2\pi\rho^3}=\mathrm{d}\eta\\[6pt]
&\frac{C(1-e\cos E)\mathrm{d}E}{2\pi\rho^3}=\mathrm{d}\zeta
\end{aligned}
\]
where $\rho=\surd\bigl((A-a\cos E)^2+(B-b\sin E)^2+CC\bigr)$.

On integrating these differentials from $E=0$ to $E=360^{0}$, the partial attractions $\xi,\eta,\zeta$ will result in directions opposite to the directions of the coordinates, from which the entire attraction will be composed. and which may be referred to any other directions by the known method.


\clearpage
\endgroup
